\documentclass[twocolumn]{aastex631}
\begin{document}
\title{The reionization ionizing-photon-budget shortfall for star-forming galaxies is not robust to systematics at z\~{}11 (optimistic systematic corner)}
\author{NebulaMind Lab (autonomous pipeline)}
\affiliation{NebulaMind Lab, \url{https://lab.nebulamind.net}}
\begin{abstract}
Reionization ionizing-photon-budget reconciliation at z\~{}11: star-forming galaxies require f\_esc=0.245 (+0.211/-0.112) to close the budget (Madau-Dickinson SFRD, log xi\_ion=25.5+/-0.15, clumping C=2-5, JWST-SFRD tail), versus indirect-proxy-inferred f\_esc=0.080 (+0.147/-0.051) from LzLCS O32/beta calibrations. Median delta(required-inferred)=+0.145 dex-frac (16-84\%: -0.013 to +0.358); 82\% of the systematic MC shows a shortfall. Conclusion: the budget CLOSES within the systematic, and the sign holds under both O32 and beta calibrations. The result is bounded by the xi\_ion x clumping x proxy-calibration systematic, not by statistics. Generated autonomously from public data (jwst) via the NebulaMind Lab runner: a bounded, reproducible, descriptive study.
\end{abstract}
\section{Introduction}
Recent studies have highlighted potential discrepancies in the ionizing photon budget during the reionization epoch, suggesting that star-forming galaxies may not produce enough photons to account for the observed reionization process [Muñoz2024]. This has led to concerns about a "photon budget crisis" and the need for a more accurate assessment of the ionizing photon production. Previous research has explored various aspects of this issue, including the role of galaxy ionizing photon budgets at lower redshifts [Duncan2015] and the challenges posed by absorption-dominated reionization scenarios [Davies2021]. However, a comprehensive analysis of the ionizing photon budget at z\~{}11 is still lacking.
\section{Data and method}
To address this gap, we employ a literature-anchored budget calculation approach that relies on published values for key parameters. Specifically, we use the cosmic star formation rate density (SFRD) from Madau \& Dickinson's [Madau2017] analytic fitting function and adopt the xi\_ion and O32/beta f\_esc proxy calibrations from Chisholm+22, Flury+22, and Simmonds+24. By combining these elements, we can estimate the required escape fraction (f\_esc) for star-forming galaxies to close the ionizing photon budget at z\~{}11.
\section{Result}
Our analysis reveals that star-forming galaxies require an escape fraction of f\_esc=0.245 (+0.211/-0.112) to reconcile the reionization ionizing-photon-budget at z\~{}11, assuming a Madau-Dickinson SFRD, log xi\_ion=25.5+/-0.15, and clumping factor C=2-5. In contrast, indirect-proxy-inferred f\_esc values from LzLCS O32/beta calibrations yield f\_esc=0.080 (+0.147/-0.051). The median difference between the required and inferred escape fractions is +0.145 dex-frac (16-84\%: -0.013 to +0.358), with 82\% of systematic Monte Carlo simulations showing a shortfall.
\begin{figure}\centering\includegraphics[width=\columnwidth]{result.png}\caption{The reionization ionizing-photon-budget shortfall for star-forming galaxies is not robust to systematics at z\~{}11 (optimistic systematic corner)}\end{figure}
\section{Caveats}
It is essential to acknowledge the limitations of our approach, which relies on automated, single-selection, uncalibrated measurements. The accuracy of our results depends heavily on the assumptions and calibrations adopted from previous studies, which may introduce biases or uncertainties. Furthermore, our analysis does not account for potential variations in galaxy properties or environmental factors that could influence the ionizing photon budget. Therefore, while our findings provide valuable insights into the reionization process, they should be interpreted with caution and considered alongside other observational and theoretical efforts to fully understand this complex phenomenon.
\section*{References}
\footnotesize
[Muoz2024] Muñoz et al. (2024). Reionization after JWST: a photon budget crisis?. bibcode 2024MNRAS.535L..37M \par [Park2022] Park et al. (2022). Calibrating excursion set reionization models to approximately conserve ionizing photons. bibcode 2022MNRAS.517..192P \par [Duncan2015] Duncan et al. (2015). Powering reionization: assessing the galaxy ionizing photon budget at z < 10. bibcode 2015MNRAS.451.2030D \par [Davies2021] Davies et al. (2021). The Predicament of Absorption-dominated Reionization: Increased Demands on Ionizing Sources. bibcode 2021ApJ...918L..35D \par [Madau2017] Madau (2017). Cosmic Reionization after Planck and before JWST: An Analytic Approach. bibcode 2017ApJ...851...50M

\end{document}
